
We opened this paper with an inefficiency: foundation model training on datasets where 30-40\% of tokens are duplicates, wasting millions in compute on data that teaches nothing new. Scaling laws optimize model size N and dataset size D, but treat all tokens as equivalent.

We close with a solution: a validated measurement framework that quantifies this waste. Our four-component Q(D) metric---deduplication ratio, domain diversity, perplexity score, and token efficiency---correlates strongly (r = 0.78, $R^2$ = 0.61) with information density across 120GB of controlled C4 subsets. Deduplication emerges as the strongest predictor (r = 0.72), validating GPT-3's empirical heuristic with quantitative evidence. All four components contribute non-redundantly with low measurement variance (CV < 10\%), meeting the stability threshold for production deployment.

This contribution is deliberately narrow: we provide a measurement tool, not a complete scaling law. We have not proven that curation causally increases information density, nor tested compute-quality tradeoff curves. These gaps are documented in the discussion.

But the tool itself is valuable. For decades, practitioners have known intuitively that ''data quality matters''---they deduplicate, filter, mix domains---yet these decisions remained heuristic. GPT-3 used fuzzy dedup because ''it seemed to help.'' The Pile emphasized diversity because ''breadth improves generalization.'' No one could quantify the effect size or prioritize investments.

Now they can measure. Deduplication (r = 0.72) provides higher return than efficiency optimization (r = 0.53). Domain diversity (r = 0.65) explains 42\% of information density variance---not dominant, but substantial. A composite Q(D) score computed before training begins predicts 61\% of data's learning value. These numbers enable data-driven curation decisions.

Foundation model scaling laws have long optimized two dimensions: model size N and dataset size D. With validated Q(D) measurement, we can now optimize the third dimension---data quality itself.
